% chapters/11/section-05.tex
\section{Menyusun Skenario Uji dan Merumuskan Perbaikan}
\label{sec:skenario-perbaikan}

Bagian ini memuat dua indikator sekaligus: menyusun skenario uji
(indikator \ind{a}) dan merumuskan perbaikan berdasarkan hasilnya
(indikator \ind{c}).

\subsection{Syarat Skenario Uji yang Berguna}

Skenario uji pada kecerdasan buatan berbeda dari pengujian program biasa.
Apabila pengujian biasa menuntut jawaban benar atau salah, skenario uji
perilaku menuntut keadaan yang memicu perilaku tertentu, lalu pengamatan atas
mutu perilaku itu.

Skenario yang berguna memenuhi empat syarat.

\begin{enumerate}
  \item \textbf{Dapat diulang dengan hasil yang sama.} Posisi pemain dan
        urutannya harus tetap, bukan bergantung pada kecakapan orang yang
        menjalankannya. Cara yang dipakai pada bab ini adalah jalur pemain yang
        ditentukan oleh nomor bingkai.
  \item \textbf{Memicu perilaku yang hendak dinilai.} Apabila yang dinilai
        adalah perilaku mengejar, skenario harus membuat pemain terlihat
        selama jangka waktu yang cukup panjang.
  \item \textbf{Memuat ukuran yang dapat dibandingkan.} Setiap skenario harus
        menghasilkan angka, misalnya jarak rata-rata, banyaknya pencarian, atau
        waktu yang diperlukan. Angka itulah yang dibandingkan antar-skenario.
  \item \textbf{Mewakili keadaan pemakaian yang sesungguhnya.} Skenario pada
        ruang terbuka saja tidak cukup; permainan yang sesungguhnya juga memuat
        lorong sempit dan tikungan.
\end{enumerate}

\subsection{Tiga Skenario yang Diminta}

RPS mata kuliah ini menuntut tiga skenario uji pada mini proyek AI. Ketiganya
dipilih karena memicu kesulitan yang berbeda, dan ketiganya dapat dibangun
dengan alat yang telah tersedia. Rinciannya pada
Tabel~\ref{tab:tiga-skenario}.

\begin{table}[htbp]
\centering
\caption{Tiga skenario uji beserta keadaan yang dipicu dan yang diamati}
\label{tab:tiga-skenario}
{\setlength{\tabcolsep}{4pt}\footnotesize%
\begin{tabularx}{\textwidth}{@{}L{2.8cm} L{4.0cm} Y@{}}
\toprule
\headrow \textbf{Skenario} & \textbf{Keadaan yang dipicu} & \textbf{Yang diamati} \\
\midrule
Ruang terbuka &
Pemain berlari lurus menjauh, tanpa rintangan di antaranya. &
Apakah NPC dapat menyusul pemain, dan berapa lama waktu yang diperlukan.
Pada ruang terbuka, jalur lurus sudah mencukupi, sehingga hasilnya menjadi
pembanding bagi kedua skenario berikutnya. \\
\addlinespace
Labirin &
Pemain berlari mengelilingi lorong, sehingga jalan lurus selalu terhalang. &
Apakah NPC menemukan jalan memutar, dan apakah jalurnya berubah-ubah tanpa
sebab. Perhatikan pula banyaknya pencarian jalur yang dijalankan. \\
\addlinespace
Rintangan dinamis &
Sebuah benda bergerak menutupi lorong yang sedang dilalui NPC. &
Apakah NPC dapat mengatasi jalur yang menjadi usang di tengah jalan, dan
berapa lama waktu yang diperlukan untuk memperbaikinya. Skenario ini
mengungkap kelemahan yang paling penting dari rancangan pada bab ini. \\
\bottomrule
\end{tabularx}}
\end{table}

\subsection{Bentuk Laporan Pengamatan}

Hasil pengamatan sebaiknya dicatat dengan bentuk yang tetap, agar dapat
dibandingkan. Bentuk yang dipakai pada mata kuliah ini memuat empat hal:
nama skenario, keadaan yang dipicu, angka yang terukur, dan kesan mutu
perilakunya.

Angka yang terukur diambil dari program, sedangkan kesan mutunya berasal dari
pengamatan Anda. Keduanya diperlukan: angka tanpa kesan tidak menjelaskan
apakah perilakunya terasa baik, sedangkan kesan tanpa angka tidak dapat
dibandingkan antar-skenario.

\subsection{Merumuskan Perbaikan}

Perbaikan yang baik memiliki tiga sifat: berdasarkan pengukuran, terarah pada
penyebabnya, dan disertai cara mengukur keberhasilannya. Tabel~\ref{tab:perbaikan}
memuat perbaikan yang penulis rumuskan untuk mini proyek ini, sebagai contoh
bentuk perumusan yang dimaksud.

\begin{table}[htbp]
\centering
\caption{Perbaikan yang dirumuskan berdasarkan hasil evaluasi}
\label{tab:perbaikan}
{\setlength{\tabcolsep}{4pt}\footnotesize%
\begin{tabularx}{\textwidth}{@{}L{3.4cm} L{4.2cm} Y@{}}
\toprule
\headrow \textbf{Temuan} & \textbf{Usulan perbaikan} & \textbf{Cara mengukur keberhasilannya} \\
\midrule
Pencarian setiap bingkai tidak membawa perbaikan perilaku &
Pakai selang hitung ulang 0,3 detik sebagai bawaan. &
Banyaknya pencarian turun menjadi sepertiga, sedangkan jarak rata-rata tidak
membesar. \\
\addlinespace
Seluruh NPC menghitung ulang pada bingkai yang sama, sehingga laju bingkai naik
turun &
Sebar waktu perhitungan dengan memberikan awal yang berbeda kepada setiap NPC. &
Simpangan laju bingkai antar-bingkai menyusut, sedangkan rata-ratanya tidak
berubah. \\
\addlinespace
Jalur dihitung ulang walaupun jalur yang lama masih memadai &
Periksa lebih dahulu apakah ujung jalur masih dekat dengan pemain; apabila ya,
perpanjang jalur itu alih-alih menghitungnya dari awal. &
Banyaknya pencarian turun tanpa membesarnya jarak rata-rata. \\
\addlinespace
Gerak NPC terlihat menyudut pada tikungan &
Singkirkan petak yang segaris, sehingga hanya titik belok yang tersisa. &
Pengamatan atas geraknya, dan penambahan waktu perhitungan yang masih dapat
diterima. \\
\bottomrule
\end{tabularx}}
\end{table}

\begin{catatabox}
Tiga perbaikan pertama menuntut pengukuran yang telah Anda
miliki, sedangkan perbaikan keempat menuntut pengamatan yang lebih halus.
Membedakan keduanya penting: perbaikan yang dapat diukur dengan angka sebaiknya
didahulukan, sebab keberhasilannya dapat dibuktikan, bukan hanya dirasakan.
\end{catatabox}
